\documentclass[11pt]{article}
\usepackage[margin=1in]{geometry}
\usepackage{amsmath,amssymb,amsthm}

\newtheorem{definition}{Definition}
\newtheorem{proposition}{Proposition}
\newtheorem{remark}{Remark}

\newcommand{\phig}{\varphi}
\newcommand{\mustar}{\mu_\star}
\newcommand{\Epass}{E_{\mathrm{pass}}}
\newcommand{\Ecoh}{E_{\mathrm{coh}}}
\newcommand{\Bpow}{B_{\mathrm{pow}}}

\title{The Recognition Science Fermion Mass Backbone at the Anchor Scale\\
\large Structural Mass Coordinates, Sector Yardsticks, and the Explicit Calibration Seam}

\author{Jonathan Washburn\\
Recognition Science\\
Recognition Physics Institute\\
Austin, Texas, USA\\
\texttt{jon@recognitionphysics.org}}

\date{March 2026}

\begin{document}
\maketitle

\begin{abstract}
This paper isolates the backbone of the Recognition Science fermion-mass framework and separates three layers that were mixed in earlier drafts: the anchor-scale structural coordinate law, the off-anchor Standard Model transport layer, and the optional SI reporting seam. The structural law is written in RS-native units as
\[
  m_f^{\mathrm{RS}}(\mustar)
  = A_{s(f)}\,\phig^{\,r_f - 8 + \mathrm{gap}(Z_f)}
\]
at the common charged-fermion anchor \(\mustar = 182.201\,\mathrm{GeV}\), where
\[
  A_s = 2^{\Bpow(s)}\Ecoh\,\phig^{r_0(s)}.
\]
The cube integers \(V=8\), \(E=12\), \(F=6\), \(A=1\), \(\Epass=11\), and \(W=17\) furnish the integer vocabulary. The canonical yardstick pairs \((\Bpow,r_0)\) are \((-22,62)\) for leptons, \((-1,35)\) for up-type quarks, \((23,-5)\) for down-type quarks, and \((1,55)\) for the electroweak sector. Generation torsion is \(\tau_g\in\{0,11,17\}\), the octave reference is \(-8\), and the charge map is encoded by the integer \(Z\). The paper does not claim that anchor-scale structural masses are pole masses, does not identify the structural band exponent with the Standard Model transport exponent, and does not hide SI conversion inside the model layer. It also states explicitly which deeper derivations remain frontier items: the charge-reversal normalization underlying the canonical gap map, the gauge-to-cube sector identification, and the lepton sub-leading corrections. The purpose is not to finish the whole spectrum in one stroke, but to provide a clean mass-law backbone to which phenomenology papers can attach without ambiguity.
\end{abstract}

\section{Introduction}

The Standard Model organizes fermion masses through Yukawa couplings but does not derive the observed pattern of charged-lepton and quark hierarchies from a smaller combinatorial substrate. The Recognition Science program proposes that the mass pattern is controlled by a discrete cost law, a ledger structure, a forced golden-ratio ladder, and the combinatorics of the three-cube \(Q_3\). The present paper has a narrower goal than a full spectrum paper. It writes down the canonical mass-law backbone, fixes its claim boundary, and makes the reporting protocol explicit.

That narrower goal matters. A mass framework becomes hard to evaluate when three different objects are mixed together: the structural coordinate law at a single common anchor, the Standard Model renormalization-group transport used to compare different scales and schemes, and the SI conversion used to display dimensionless RS-native outputs as MeV or GeV. These are different layers. They should be written separately, audited separately, and falsified separately.

The backbone presented here is therefore a methods paper. It gives the common coordinate law, the yardstick table, the charge map, the generation-torsion rule, the anchor protocol, and the explicit calibration seam. Companion papers can then focus on charged leptons, CKM observables, neutrinos, or later quark refinements without re-litigating the same bookkeeping on every page.

\section{Scope and status conventions}

Three status words are used throughout this paper.

\paragraph{Structural.} A structural statement belongs to the fixed model layer at the anchor. It is part of the canonical mass coordinate system and is used before any SI display or off-anchor transport.

\paragraph{Protocol.} A protocol statement is an explicit convention required for comparison, such as the common anchor scale, the transport policy, or the single-scalar SI seam. Protocol statements are not hidden fit parameters; they are declared global rules.

\paragraph{Frontier.} A frontier statement is not claimed here as a closed first-principles derivation, even if a canonical formula is already in use in the model layer. Frontier items are kept visible rather than hidden behind suggestive prose.

This paper makes five concrete claims.

\begin{enumerate}
  \item The mass model at the anchor factorizes into a sector yardstick and a species exponent.
  \item The integer vocabulary is fixed by the three-cube: \(V=8\), \(E=12\), \(F=6\), \(A=1\), \(\Epass=11\), and \(W=17\).
  \item Once the sector labels are fixed, the yardstick integers \((\Bpow,r_0)\) and the generation-torsion data are no longer opaque literals; they are computed from the counting layer.
  \item Off-anchor comparison is handled by a separate Standard Model transport exponent \(f^{\mathrm{RG}}\), which is not the same object as the structural band exponent \(\mathrm{gap}(Z)\).
  \item SI reporting requires an explicit one-scalar seam. Species-by-species calibration is forbidden.
\end{enumerate}

This paper does \emph{not} claim a closed derivation of the charge-reversal normalization used in the canonical gap map, a closed first-principles derivation of the gauge-sector-to-cube-role identification, or a closed derivation of the sub-leading charged-lepton corrections. Those items are part of the frontier, and they are kept explicit.

\section{From the cost law to the cube integers}

The starting point is the Recognition Composition Law,
\[
  J(xy) + J(x/y) = 2J(x)J(y) + 2J(x) + 2J(y),
  \qquad x,y>0,
\]
together with the normalization \(J(1)=0\) and the local curvature calibration at the minimum. The unique positive cost function compatible with these conditions is
\[
  J(x) = \frac{1}{2}\left(x + \frac{1}{x}\right) - 1.
\]

Self-similar discrete closure forces the ladder base \(\phig\) as the positive solution of
\[
  x^2 = x + 1,
  \qquad
  \phig = \frac{1+\sqrt{5}}{2}.
\]
A ledger-compatible walk then forces an eight-step minimal cycle and singles out three spatial dimensions. The counting layer is therefore the three-cube \(Q_3\), and its canonical integers are
\[
  V=8,
  \qquad
  E=12,
  \qquad
  F=6,
  \qquad
  A=1,
  \qquad
  \Epass = E-A = 11,
  \qquad
  W = \Epass + F = 17.
\]

These six integers are the entire small-integer vocabulary needed in the mass backbone. No particle species gets its own bespoke integer inventory.

\section{The anchor-scale mass law}

\begin{definition}[Anchor-scale structural mass law]
The canonical RS-native mass coordinate of a fermion species \(f\) at the common charged-fermion anchor \(\mustar\) is
\[
  m_f^{\mathrm{RS}}(\mustar)
  = A_{s(f)}\,\phig^{\,r_f - 8 + \mathrm{gap}(Z_f)}.
\]
Here \(s(f)\) is the sector label, \(A_s\) is the sector yardstick, \(r_f\) is the rung, \(-8\) is the octave reference, and \(Z_f\) is the integerized charge coordinate.
\end{definition}

The anchor used for the charged-fermion display is
\[
  \mustar = 182.201\,\mathrm{GeV}.
\]
Its role is purely coordinatizing: it is the single common scale at which the structural bands are stated. This is not a claim that the anchor-scale structural masses are identical to pole masses or to \(\overline{\mathrm{MS}}\) masses. Comparison to those observables is a second step and uses the transport layer introduced later.

The sector yardstick has the factorized form
\[
  A_s = 2^{\Bpow(s)}\Ecoh\,\phig^{r_0(s)},
  \qquad
  \Ecoh = \phig^{-5}.
\]
The canonical band map is
\[
  \mathrm{gap}(Z)
  = \log_{\phig}\!\left(1+\frac{Z}{\phig}\right)
  = \frac{\ln\!\left(1+Z/\phig\right)}{\ln\phig}.
\]
The exponent \(-8\) is the octave reference associated with one full eight-tick closure. It is part of the anchor display itself.

\begin{proposition}[Equal-\(Z\) family ratios]
If two species \(f\) and \(g\) lie in the same sector and carry the same integerized charge coordinate, then
\[
  \frac{m_f^{\mathrm{RS}}(\mustar)}{m_g^{\mathrm{RS}}(\mustar)} = \phig^{r_f-r_g}.
\]
\end{proposition}

\begin{proof}
At fixed sector and fixed \(Z\), the factors \(A_s\), \(-8\), and \(\mathrm{gap}(Z)\) cancel in the ratio. Only the rung difference survives.
\end{proof}

This simple cancellation is one of the main reasons to separate the structural anchor law from the transport and calibration layers. It shows which predictions are genuinely ladder statements and which are not.

\section{Sector yardsticks, charge map, and rungs}

The canonical sector yardsticks are listed in the following table.

\begin{center}
\begin{tabular}{|l|c|c|c|}
\hline
Sector & \(\Bpow\) & \(r_0\) & Yardstick \\
\hline
Lepton & \(-22\) & \(62\) & \(2^{-22}\Ecoh\,\phig^{62}\) \\
Up-type quark & \(-1\) & \(35\) & \(2^{-1}\Ecoh\,\phig^{35}\) \\
Down-type quark & \(23\) & \(-5\) & \(2^{23}\Ecoh\,\phig^{-5}\) \\
Electroweak & \(1\) & \(55\) & \(2\Ecoh\,\phig^{55}\) \\
\hline
\end{tabular}
\end{center}

Once the sector labels are fixed, these pairs are computed from the cube-counting layer rather than inserted as opaque constants. The arithmetic is therefore part of the backbone. What is \emph{not} claimed here is the deeper first-principles reason that these four physical sector labels must attach to these four cube roles. That label-to-role identification remains a frontier item and is not hidden.

Generation torsion is encoded by the three-element set
\[
  \tau_g \in \{0,\Epass,W\} = \{0,11,17\}.
\]
The canonical rung rule is
\[
  r_f = \ell_{s(f)} + \tau_{g(f)},
\]
where \(\ell_s\) is the baseline rung of the sector. For the charged fermions the canonical integer rungs are
\[
  (e,\mu,\tau) = (2,13,19),
\]
\[
  (u,c,t) = (4,15,21),
  \qquad
  (d,s,b) = (4,15,21).
\]

The charge coordinate is integerized by
\[
  \tilde Q = 6Q \in \mathbb{Z}.
\]
For leptons and quarks the canonical maps are
\[
  Z_{\ell} = \tilde Q^2 + \tilde Q^4,
  \qquad
  Z_q = 4 + \tilde Q^2 + \tilde Q^4,
  \qquad
  Z_{\nu}=0.
\]
Hence the family values used in the anchor law are
\[
  Z_{\ell}=1332,
  \qquad
  Z_u=276,
  \qquad
  Z_d=24,
  \qquad
  Z_{\nu}=0.
\]
The corresponding canonical band exponents are
\[
  \mathrm{gap}(1332)\approx 13.953,
  \qquad
  \mathrm{gap}(276)\approx 10.692,
  \qquad
  \mathrm{gap}(24)\approx 5.740,
  \qquad
  \mathrm{gap}(0)=0.
\]

These are large structural band coordinates at the anchor. They are not small perturbative residues.

\section{Transport is not structure}

The most important bookkeeping rule in the mass framework is that the structural band exponent and the Standard Model transport exponent are different quantities.

\begin{definition}[Transported mass display]
To compare the anchor-scale structural law to a target scale \(\mu\), one adds a separate transport exponent:
\[
  m_f^{\mathrm{RS}}(\mu)
  = A_{s(f)}\,\phig^{\,r_f - 8 + \mathrm{gap}(Z_f) + f_f^{\mathrm{RG}}(\mustar,\mu)}.
\]
The quantity \(f_f^{\mathrm{RG}}\) is the Standard Model transport exponent associated with a fixed renormalization-group policy.
\end{definition}

The transport exponent is small. The gap exponent is typically of order one to ten. Treating them as the same object is a category mistake. The gap is a structural band coordinate stated at the common anchor. The transport term is a conventional scheme-and-scale correction used only when one moves away from that anchor.

This distinction also prevents a familiar cheat. One must not absorb threshold choices, running-coupling choices, or scale choices into the structural law and then pretend they were part of the combinatorics all along. The structural layer ends with \(A_s\), \(r_f\), \(-8\), and \(\mathrm{gap}(Z_f)\). Everything else belongs to transport.

\section{The explicit SI calibration seam}

RS-native masses are dimensionless coordinates in the native mass-energy gauge. A paper written for particle phenomenology will usually display masses as energies in eV, MeV, or GeV. That requires a declared calibration seam.

Let \(\tau_0\) denote the number of SI seconds per atomic tick. Then the two reporting factors are
\[
  \ell_0^{\mathrm{SI}} = c\tau_0,
  \qquad
  \epsilon_{\mathrm{coh}}^{\mathrm{SI}} = \frac{\hbar}{\tau_0}.
\]
In the particle-mass context one normally reports \(mc^2\) in energy units, so the displayed energy corresponding to an RS-native mass coordinate is obtained by multiplying by \(\epsilon_{\mathrm{coh}}^{\mathrm{SI}}\). The crucial point is that \(\tau_0\) is a \emph{single global scalar}. It is not chosen separately for each species.

\begin{remark}[Forbidden moves]
The following moves are outside the backbone and would count as hidden fitting.
\begin{enumerate}
  \item Choosing a separate calibration for each species.
  \item Using a measured mass to set that same species' scale.
  \item Hiding transport-policy choices inside the structural exponent.
\end{enumerate}
\end{remark}

Making the seam explicit does not weaken the model layer. It strengthens it, because it forbids unit-smuggling. The model layer is responsible for the dimensionless structure. The seam is responsible for SI display. Mixing the two only creates fake precision.

\section{Why the anchor is not a per-species fit}

The natural criticism is that an anchor selected after looking at the spectrum could be circular. The backbone avoids that criticism only if two conditions are respected.

First, the anchor must be \emph{common}. It is one scale for the charged-fermion display, not a separate scale for each family or each particle.

Second, the numerical value of the anchor must be obtained from a fixed, species-independent stationarity protocol in the Standard Model transport layer. Once that policy is fixed, the anchor is part of the global display protocol, not a knob that can be adjusted particle by particle.

Under those conditions the anchor plays the same role for every charged fermion. It is a shared coordinate origin for the structural bands. It is not a disguised Yukawa fit. The moment one allows family-specific anchors, the backbone collapses into bookkeeping theater.

\section{Frontier items kept explicit}

A publishable backbone paper should say plainly what is closed and what is not. The following items remain frontier statements in the present manuscript.

\paragraph{Gap normalization.} The canonical band map \(\mathrm{gap}(Z)=\log_{\phig}(1+Z/\phig)\) is used throughout the model layer, but this paper does not claim a closed first-principles derivation of the charge-reversal normalization that singles out the full canonical calibration.

\paragraph{Sector identification.} The table of yardstick integers is fixed once the sector labels are chosen, but this paper does not claim a closed first-principles derivation of why the physical gauge sectors must attach to those four cube roles and no others.

\paragraph{Charged-lepton sub-leading terms.} The backbone paper does not derive the sub-leading charged-lepton corrections involving \(1/(4\pi)\), \(\alpha^2\), and \((2W+3)\alpha/2\). Those belong in the charged-lepton phenomenology paper and should be kept separate from the backbone.

\paragraph{Quark refinement layer.} The present paper uses the core integer-rung anchor display. It does not claim an equivalence between that core display and any quarter-ladder quark scaffold used elsewhere. Any such equivalence must be proved separately or left explicitly conditional.

These frontier statements do not invalidate the backbone. They simply mark where the current backbone ends and where later papers begin.

\section{Falsifiers and mandatory ablations}

A backbone that cannot be broken is not science. The present one has clear failure modes.

\begin{enumerate}
  \item If no single common anchor can support a stable charged-fermion display under a fixed transport policy, the backbone fails.
  \item If same-sector equal-\(Z\) mass ratios do not reduce to \(\phig\)-powers of rung differences, the rung law fails.
  \item If SI display requires more than one calibration scalar, the explicit seam fails and the model has hidden parameters.
  \item If a simpler charge map reproduces the same structural bands with fewer declared ingredients, then the present charge integerization is not uniquely justified.
\end{enumerate}

Ablation tests are equally important. Any companion phenomenology paper should compare its preferred formula against the skeleton law
\[
  m_{\mathrm{skel}}(\mustar)=A_s\,\phig^{r-8},
\]
which removes the charge band entirely. If the skeleton law were already sufficient, then the \(Z\)-map would be decorative rather than structural. A serious paper must show that the added structure earns its keep.

Likewise, any paper that uses off-anchor data must show that its conclusions are stable under a fixed transport policy and do not depend on quietly changing loop order, thresholds, or matching conventions from one family to another.

\section{Conclusion}

The Recognition Science mass framework becomes much easier to judge once its layers are separated cleanly. The backbone consists of the three-cube counting layer, the golden-ratio ladder, the sector yardsticks, the rung rule, the charge map, the octave reference, and the common anchor-scale display
\[
  m_f^{\mathrm{RS}}(\mustar)=A_{s(f)}\,\phig^{\,r_f-8+\mathrm{gap}(Z_f)}.
\]
Off-anchor comparison belongs to a separate transport layer, and SI numerals belong to a separate calibration seam.

That separation is the main result of this paper. It removes the two easiest reviewer attacks: hidden species-by-species fitting and unit-smuggling. It also creates a clean division of labor for the paper series. A charged-lepton paper can now focus on the lepton chain. A CKM paper can focus on mixing observables. A neutrino paper can focus on ladder predictions and dimensionless splitting ratios. A later quark paper can address its remaining frontier items without contaminating the backbone.

The universe is already weird enough. The bookkeeping does not need to be weird too.

\end{document}
